% ============================================================
%  CHAPTER 2 -- Fundamentals of Probability
%  Presenter 2
% ============================================================
\section{Fundamentals of Probability}

\begin{frame}{Sample Space and Events}
\begin{itemize}
    \item A \textbf{random experiment} is a process whose outcome is uncertain.
    \item The \textbf{sample space} $\Omega$ is the set of all possible outcomes.
    \item An \textbf{event} $A$ is a subset of the sample space: $A \subseteq \Omega$.
\end{itemize}
\begin{exampleblock}{Examples}
\begin{itemize}
    \item Coin toss: $\Omega = \{H, T\}$; event $A=\{H\}$.
    \item Rolling a die: $\Omega = \{1,2,3,4,5,6\}$; event ``even'' $=\{2,4,6\}$.
    \item ML: predicting a class label $\Omega = \{c_1, c_2, \dots, c_k\}$.
\end{itemize}
\end{exampleblock}
\end{frame}

\begin{frame}{Axioms of Probability}
\textbf{Kolmogorov Axioms:}
\begin{enumerate}
    \item \textbf{Non-negativity:} $P(A) \geq 0$ for every event $A$.
    \item \textbf{Normalization:} $P(\Omega) = 1$.
    \item \textbf{Additivity:} For mutually exclusive events $A_1, A_2, \dots$:
    \[
        P\!\left(\bigcup_{i=1}^{\infty} A_i\right) = \sum_{i=1}^{\infty} P(A_i)
    \]
\end{enumerate}
\vspace{4pt}
\textbf{Consequences:}
\begin{itemize}
    \item $P(\emptyset) = 0$
    \item $P(A^c) = 1 - P(A)$
    \item $0 \le P(A) \le 1$
\end{itemize}
\end{frame}

\begin{frame}{Conditional Probability}
\begin{block}{Definition}
\[
P(A \mid B) = \frac{P(A \cap B)}{P(B)}, \quad P(B) > 0
\]
\end{block}
\begin{itemize}
    \item The probability of event $A$ occurring \emph{given} that $B$ has occurred.
    \item Fundamental to updating beliefs with new evidence.
    \item \textbf{Example:} $P(\text{spam} \mid \text{``free''})$ -- probability an email is spam given it contains the word ``free''.
\end{itemize}
\end{frame}

\begin{frame}{Independence of Events}
Two events $A$ and $B$ are \textbf{independent} if:
\[
P(A \cap B) = P(A) \cdot P(B)
\]
Equivalently: $P(A \mid B) = P(A)$ -- knowing $B$ tells us nothing about $A$.
\vspace{6pt}
\begin{itemize}
    \item \textbf{Naive Bayes} assumes all features are conditionally independent given the class.
    \item True independence is rare in practice, but the assumption often yields surprisingly good results.
    \item \textbf{Mutual independence} of $n$ events requires all $2^n - n - 1$ subset conditions to hold.
\end{itemize}
\end{frame}

\begin{frame}{Law of Total Probability}
If $B_1, B_2, \dots, B_n$ form a \textbf{partition} of $\Omega$ (mutually exclusive, exhaustive):
\[
\boxed{P(A) = \sum_{i=1}^{n} P(A \mid B_i)\, P(B_i)}
\]
\begin{itemize}
    \item Breaks a complex probability into simpler conditional pieces.
    \item Essential building block for \textbf{Bayes' theorem}.
    \item \textbf{ML application:} Marginalising over latent variables:
    \[
        P(x) = \sum_{z} P(x \mid z)\, P(z)
    \]
\end{itemize}
\end{frame}

\begin{frame}{Numerical Example: Conditional Probability}
\begin{exampleblock}{Problem}
A bag has 5 red and 3 blue balls. Two balls are drawn without replacement.\\
Find $P(\text{2nd is red} \mid \text{1st is red})$.
\end{exampleblock}
\textbf{Solution:}
\begin{itemize}
    \item After drawing 1 red ball: remaining = 4 red + 3 blue = 7 balls.
    \item $P(\text{2nd red} \mid \text{1st red}) = \dfrac{4}{7} \approx \mathbf{0.571}$
\end{itemize}
\textbf{Verification using formula:}
\[
P(A \mid B) = \frac{P(A \cap B)}{P(B)} = \frac{5/8 \times 4/7}{5/8} = \frac{4}{7} \approx 0.571 \;\checkmark
\]
\end{frame}

\begin{frame}{Numerical Example: Law of Total Probability}
\begin{exampleblock}{Problem}
Factory has 3 machines. Machine A produces 50\% of items (2\% defective), Machine B produces 30\% (3\% defective), Machine C produces 20\% (5\% defective).\\
Find the probability that a randomly chosen item is defective.
\end{exampleblock}
\textbf{Solution:}
\begin{align*}
P(D) &= P(D\mid A)\,P(A) + P(D\mid B)\,P(B) + P(D\mid C)\,P(C)\\
     &= (0.02)(0.50) + (0.03)(0.30) + (0.05)(0.20)\\
     &= 0.010 + 0.009 + 0.010\\
     &= \mathbf{0.029} \quad \text{(2.9\% defective overall)}
\end{align*}
\end{frame}
